% Figure 3. The whole model in one picture: what a save writes, what a version
% is, and where the bytes in your files come from. Same session as Figure 1.
\newcommand{\FigThreeCaption}{What one save writes down, what a version is, and
where the bytes in a file come from. A save writes
one record per changed function and files the records under a guess at the request
they served, which the developer corrects with one command (left). A version is
the set of those records, so taking out the price cache is subtracting six records
from eighteen, and two rules decide which sets are allowed (middle). The files are
rebuilt from the set with the text between functions kept exactly, so the same set
always produces the same bytes (right). Two kinds of mark in the right-hand column
are the two consequences \sgt{} reports about this revert and does not act on. The hollow
triangle on \protect\fsym{api.py::handle} means the retry policy edited that
function after the cache did, so \sgt{} took the cache's lines back out of the
current text and the developer's test suite decides whether the result stands. The
open circles are three functions that still call
\protect\fsym{cache.py::get\_cached}, which \sgt{} names and leaves alone.
\protect\fsym{db.py::query} is the one function that goes back to text it had
earlier in the day.}

\newcommand{\FigThreeDescription}{Three panels. The left panel shows a terminal
transcript of one save command, followed by the three feature groups the save
landed in, one of them newly created and unnamed, and then a card listing the
five fields one edit record holds. The middle panel lists the fourteen functions
this session changed, grouped by file, with the edit marks for each on a short
rail, in two columns. The left column is the set as it stands and the right
column is the same list after the price cache has been reverted, where the cache
edits are drawn as dotted outlines, one file is marked as removed, one function is
marked with a hollow triangle, and three functions are marked with an open circle.
The right panel draws the file
api.py twice as a vertical stack of bands, once before and once after, where
coloured bands are function bodies and grey bands are the text between them,
which is identical in both.}

\begin{tikzpicture}[x=1mm, y=1mm, line cap=round, font=\sffamily]

\node[hd, anchor=west] at (0,78)   {1\quad WHAT A SAVE WRITES};
\node[hd, anchor=west] at (52,78)  {2\quad WHAT A VERSION IS};
\node[hd, anchor=west] at (117,78) {3\quad WHAT IS IN YOUR FILES};
\draw[color=sgtGrey!40, line width=0.3pt] (0,76.5) -- (176,76.5);
\draw[color=sgtGrey!30, line width=0.3pt] (49.5,11) -- (49.5,76.5);
\draw[color=sgtGrey!30, line width=0.3pt] (114,11) -- (114,76.5);
\draw[color=sgtGrey!40, line width=0.3pt] (0,9.6) -- (176,9.6);

% ==================== panel 1: what a save writes =========================
\node[anchor=north west, align=left, inner sep=0, text=sgtInk,
      font=\ttfamily\fontsize{5.2}{7}\selectfont] at (0,74.6)
  {\$ sgt save -m "retry the upstream call\\
   \hphantom{\$ sgt save -m "}when it returns a 5xx"\\
   \ \ save 7f10de\ \ 6 edits};

\oprew{1.3}{62.2}{fRetry}
\node[anchor=north west, align=left, inner sep=0, text=sgtInk,
      font=\ttfamily\fontsize{5.2}{6.6}\selectfont] at (3.4,63.2)
  {\textcolor{fRetry}{new feature} (7c02b4d1)\ \ 3 edits\\
   retry.py::Retry\ \ retry.py::\_backoff\\
   api.py::handle\\
   unnamed. name it with\\
   sgt feature rename 7c02b4d1 "retry policy"};

\oprew{1.3}{45.6}{fCache}
\node[anchor=north west, align=left, inner sep=0, text=sgtInk,
      font=\ttfamily\fontsize{5.2}{6.6}\selectfont] at (3.4,46.6)
  {\textcolor{fCache}{price cache} (001f73c2)\ \ 2 edits\\
   cache.py::get\_cached\ \ cache.py::\_ttl};

\oprew{1.3}{39.4}{sgtGrey}
\node[anchor=north west, align=left, inner sep=0, text=sgtInk,
      font=\ttfamily\fontsize{5.2}{6.6}\selectfont] at (3.4,40.4)
  {\textcolor{sgtGrey}{row accessors} (9d20b18a)\ \ 1 edit\\
   util.py::audit};

% one record, in full
\draw[card, color=sgtGrey!45, fill=sgtPaper] (0,14.6) rectangle (47,34.2);
\fill[fRetry] (0,31.6) rectangle (47,34.2);
\node[anchor=west, inner sep=0, text=sgtPaper,
      font=\ttfamily\fontsize{5.2}{6}\selectfont] at (1.2,32.9) {edit 4b1c9f28};
\node[anchor=north west, align=left, inner sep=0, text=sgtInk,
      font=\ttfamily\fontsize{5.2}{6.9}\selectfont] at (1.2,30.6)
  {function\ \ \ \ api.py::handle\\
   kind\ \ \ \ \ \ \ \ extended\\
   started at\ \ 9ac7e1\\
   produced\ \ \ \ b41c08\\
   built on\ \ \ \ retry.py::Retry};

% ==================== panel 2: what a version is ==========================
\node[anchor=north west, align=left, text width=61mm, inner sep=0, text=sgtInk!85,
      font=\sffamily\fontsize{5.6}{6.6}\selectfont] at (52,74.8)
  {\textbf{rule 1}\ \ no edit without the edits it was built on.\ \ \ \
   \textbf{rule 2}\ \ one live version of each function.};

\node[anchor=north west, align=left, inner sep=0, text=sgtInk,
      font=\sffamily\bfseries\fontsize{5.6}{6.4}\selectfont] at (52,68.6)
  {THE SET NOW};
\node[anchor=north west, inner sep=0, text=sgtGrey,
      font=\sffamily\fontsize{5.2}{6}\selectfont] at (52,66.0)
  {18 edits, 14 functions};
\node[anchor=north west, align=left, inner sep=0, text=sgtInk,
      font=\ttfamily\fontsize{5.4}{6.2}\selectfont] at (84,68.6)
  {\$ sgt revert "price cache"};
\node[anchor=north west, inner sep=0, text=sgtGrey,
      font=\sffamily\fontsize{5.2}{6}\selectfont] at (84,66.0)
  {12 edits, 12 functions};

% #1 x, #2 y, #3 file name
\newcommand{\FTfile}[3]{\node[anchor=north west, inner sep=0, text=sgtInk!65,
  font=\sffamily\bfseries\fontsize{5.2}{6}\selectfont] at (#1,#2) {#3};}
% #1 x, #2 y, #3 symbol name, #4 colour
\newcommand{\FTsym}[4]{\node[anchor=north west, inner sep=0, text=#4,
  font=\ttfamily\fontsize{5.2}{6}\selectfont] at (#1,#2) {#3};}
% a hollow up-triangle: this version was re-drafted, and the tests decide
\newcommand{\FTredraft}[2]{\draw[sgtInk!85, line width=0.4pt, fill=sgtPaper]
  ({#1-0.95},{#2-0.85}) -- ({#1+0.95},{#2-0.85}) -- (#1,{#2+1.0}) -- cycle;}
% an open circle: this function still calls something the revert removed
\newcommand{\FTcalls}[2]{\draw[sgtInk!85, line width=0.4pt, fill=sgtPaper]
  (#1,#2) circle (0.85);}
% #1 x of rail start, #2 y of rail
\newcommand{\FTrail}[3]{\draw[rail=#3] (#1,#2) -- ({#1+9},#2);}
\newcommand{\FTraildim}[2]{\draw[raildim] (#1,#2) -- ({#1+9},#2);}

% ---- left column: the set as it stands
\FTfile{52}{63.6}{api.py}
\FTsym{54}{60.8}{handle}{sgtInk}
  \FTrail{69.5}{59.9}{fRetry}
  \opext{70.6}{59.9}{fRate}\opext{74.0}{59.9}{fCache}\opext{77.4}{59.9}{fRetry}
\FTsym{54}{58.0}{\_bucket}{sgtInk}
  \FTrail{69.5}{57.1}{fRate}\opadd{70.6}{57.1}{fRate}
\FTsym{54}{55.2}{\_key}{sgtInk}
  \FTrail{69.5}{54.3}{fRate}\opadd{70.6}{54.3}{fRate}
\FTsym{54}{52.4}{\_window}{sgtInk}
  \FTrail{69.5}{51.5}{fRate}\opadd{70.6}{51.5}{fRate}
\FTsym{54}{49.6}{export\_row}{sgtInk}
  \FTrail{69.5}{48.7}{sgtGrey}\opext{70.6}{48.7}{sgtGrey}
\FTfile{52}{46.8}{cache.py}
\FTsym{54}{44.0}{get\_cached}{sgtInk}
  \FTrail{69.5}{43.1}{fCache}
  \opadd{70.6}{43.1}{fCache}\oprew{74.0}{43.1}{fCache}
\FTsym{54}{41.2}{\_ttl}{sgtInk}
  \FTrail{69.5}{40.3}{fCache}
  \opadd{70.6}{40.3}{fCache}\opext{74.0}{40.3}{fCache}
\FTfile{52}{38.4}{db.py}
\FTsym{54}{35.6}{query}{sgtInk}
  \FTrail{69.5}{34.7}{fCache}\opext{70.6}{34.7}{fCache}
\FTsym{54}{32.8}{fetch\_row}{sgtInk}
  \FTrail{69.5}{31.9}{sgtGrey}\opmove{70.6}{31.9}{sgtGrey}
\FTfile{52}{30.0}{retry.py}
\FTsym{54}{27.2}{Retry}{sgtInk}
  \FTrail{69.5}{26.3}{fRetry}\opadd{70.6}{26.3}{fRetry}
\FTsym{54}{24.4}{\_backoff}{sgtInk}
  \FTrail{69.5}{23.5}{fRetry}\opadd{70.6}{23.5}{fRetry}
\FTfile{52}{21.6}{util.py}
\FTsym{54}{18.8}{load\_row}{sgtInk}
  \FTrail{69.5}{17.9}{sgtGrey}\opext{70.6}{17.9}{sgtGrey}
\FTsym{54}{16.0}{save\_row}{sgtInk}
  \FTrail{69.5}{15.1}{sgtGrey}\opext{70.6}{15.1}{sgtGrey}
\FTsym{54}{13.2}{audit}{sgtInk}
  \FTrail{69.5}{12.3}{sgtGrey}\opext{70.6}{12.3}{sgtGrey}

% ---- right column: the same list, six edits subtracted
\FTfile{84}{63.6}{api.py}
\FTsym{86}{60.8}{handle}{sgtInk}
  \FTrail{101.5}{59.9}{fRetry}
  \opext{102.6}{59.9}{fRate}\opghost{106.0}{59.9}{fCache}\opext{109.4}{59.9}{fRetry}
  \FTredraft{112.6}{59.9}
\FTsym{86}{58.0}{\_bucket}{sgtInk}
  \FTrail{101.5}{57.1}{fRate}\opadd{102.6}{57.1}{fRate}
\FTsym{86}{55.2}{\_key}{sgtInk}
  \FTrail{101.5}{54.3}{fRate}\opadd{102.6}{54.3}{fRate}
\FTsym{86}{52.4}{\_window}{sgtInk}
  \FTrail{101.5}{51.5}{fRate}\opadd{102.6}{51.5}{fRate}
\FTsym{86}{49.6}{export\_row}{sgtInk}
  \FTrail{101.5}{48.7}{sgtGrey}\opext{102.6}{48.7}{sgtGrey}
  \FTcalls{112.6}{48.7}
\FTfile{84}{46.8}{cache.py}
\node[anchor=east, inner sep=0, text=sgtGrey,
      font=\sffamily\itshape\fontsize{5}{5.8}\selectfont] at (113.5,46.0)
  {file removed};
\FTsym{86}{44.0}{get\_cached}{sgtGrey!80}
  \FTraildim{101.5}{43.1}
  \opghost{102.6}{43.1}{fCache}\opghost{106.0}{43.1}{fCache}
\FTsym{86}{41.2}{\_ttl}{sgtGrey!80}
  \FTraildim{101.5}{40.3}
  \opghost{102.6}{40.3}{fCache}\opghost{106.0}{40.3}{fCache}
\FTfile{84}{38.4}{db.py}
\FTsym{86}{35.6}{query}{sgtInk}
  \draw[raildim] (101.5,34.7) -- (104.0,34.7);\opghost{102.6}{34.7}{fCache}
  \node[anchor=east, inner sep=0, text=sgtGrey,
        font=\sffamily\itshape\fontsize{5}{5.8}\selectfont] at (113.5,34.7)
    {back to 14:02};
\FTsym{86}{32.8}{fetch\_row}{sgtInk}
  \FTrail{101.5}{31.9}{sgtGrey}\opmove{102.6}{31.9}{sgtGrey}
\FTfile{84}{30.0}{retry.py}
\FTsym{86}{27.2}{Retry}{sgtInk}
  \FTrail{101.5}{26.3}{fRetry}\opadd{102.6}{26.3}{fRetry}
\FTsym{86}{24.4}{\_backoff}{sgtInk}
  \FTrail{101.5}{23.5}{fRetry}\opadd{102.6}{23.5}{fRetry}
\FTfile{84}{21.6}{util.py}
\FTsym{86}{18.8}{load\_row}{sgtInk}
  \FTrail{101.5}{17.9}{sgtGrey}\opext{102.6}{17.9}{sgtGrey}
  \FTcalls{112.6}{17.9}
\FTsym{86}{16.0}{save\_row}{sgtInk}
  \FTrail{101.5}{15.1}{sgtGrey}\opext{102.6}{15.1}{sgtGrey}
  \FTcalls{112.6}{15.1}
\FTsym{86}{13.2}{audit}{sgtInk}
  \FTrail{101.5}{12.3}{sgtGrey}\opext{102.6}{12.3}{sgtGrey}

% ==================== panel 3: what is in your files ======================
\node[anchor=north west, text width=59mm, align=left, inner sep=0, text=sgtInk!85,
      font=\sffamily\fontsize{5.6}{6.6}\selectfont] at (117,74.8)
  {\sgt{} writes a file by putting each live version back in place and keeping
   every character between them.};

% #1 x left, #2 y bottom, #3 height, #4 colour, #5 label
\newcommand{\FTband}[5]{%
  \fill[#4!30, rounded corners=0.4mm] ({#1+0.6},#2) rectangle ({#1+25.4},{#2+#3});
  \node[anchor=west, inner sep=0, text=sgtInk!85,
        font=\ttfamily\fontsize{5}{5.8}\selectfont] at ({#1+2.2},{#2+#3/2}) {#5};}
% #1 x left, #2 y bottom, #3 height: text sgt keeps verbatim
\newcommand{\FTgap}[3]{%
  \fill[sgtGrey!12] ({#1+0.6},#2) rectangle ({#1+25.4},{#2+#3});
  \foreach \i in {0,1,2} {\fill[sgtGrey!45] ({#1+2.2},{#2+#3-1.1-\i*1.1})
     rectangle ({#1+2.2+13-\i*3.2},{#2+#3-0.65-\i*1.1});}}

\node[anchor=south west, inner sep=0, text=sgtInk,
      font=\ttfamily\fontsize{5.4}{6}\selectfont] at (117,68.6) {api.py};
\node[anchor=south west, inner sep=0, text=sgtInk,
      font=\ttfamily\fontsize{5.4}{6}\selectfont] at (148,68.6) {api.py};
\node[anchor=south east, inner sep=0, text=sgtGrey,
      font=\sffamily\fontsize{5}{5.8}\selectfont] at (143,68.6) {before};
\node[anchor=south east, inner sep=0, text=sgtGrey,
      font=\sffamily\fontsize{5}{5.8}\selectfont] at (174,68.6) {after};

\foreach \x in {117,148} {
  \draw[card, color=sgtGrey!45, fill=sgtPaper] (\x,20.5) rectangle ({\x+26},67.5);
  \FTgap{\x}{63.6}{3.4}
  \FTband{\x}{58.4}{4.6}{fRate}{\_bucket}
  \FTband{\x}{53.4}{4.6}{fRate}{\_key}
  \FTband{\x}{48.4}{4.6}{fRate}{\_window}
  \FTgap{\x}{44.6}{3.4}
  \FTgap{\x}{30.0}{3.4}
  \FTband{\x}{21.4}{8.2}{sgtGrey}{export\_row}
}
% handle: one version, with a bar saying which requests contributed to it
\FTband{117}{34.0}{10.2}{fRetry}{handle}
\fill[fRate]  (117.6,34.0) rectangle (118.7,37.4);
\fill[fCache] (117.6,37.4) rectangle (118.7,40.8);
\fill[fRetry] (117.6,40.8) rectangle (118.7,44.2);
\FTband{148}{34.0}{10.2}{fRetry}{handle}
\fill[fRate]  (148.6,34.0) rectangle (149.7,39.1);
\fill[fRetry] (148.6,39.1) rectangle (149.7,44.2);
\FTredraft{171.6}{39.1}

\node[anchor=north west, text width=59mm, align=left, inner sep=0, text=sgtInk!85,
      font=\sffamily\fontsize{5.6}{6.6}\selectfont] at (117,19.0)
  {The grey bands are the imports, the comments, and the blank lines, which \sgt{}
   keeps and never rewrites. The bar left of \fsym{handle} names the requests that
   contributed to the version in the file.};

% ==================== bottom: what each one buys ==========================
\node[anchor=north west, text width=47mm, align=left, inner sep=0, text=sgtInk,
      font=\sffamily\fontsize{5.8}{6.8}\selectfont] at (0,8.4)
  {Every field was read out of the files, and the sentence at the top is all the
   developer typed. No diff was sorted by hand at any point.};
\node[anchor=north west, text width=61mm, align=left, inner sep=0, text=sgtInk,
      font=\sffamily\fontsize{5.8}{6.8}\selectfont] at (52,8.4)
  {Each of \cmd{revert}, \cmd{restore}, \cmd{diff} and \cmd{switch} is one
   subtraction or one union over this list, so \sgt{} can show the resulting files
   before writing any of them.};
\node[anchor=north west, text width=59mm, align=left, inner sep=0, text=sgtInk,
      font=\sffamily\fontsize{5.8}{6.8}\selectfont] at (117,8.4)
  {A file is a rebuild from a set. Revert the cache twice and restore it once and
   the bytes match the file that was there before any of it.};

\end{tikzpicture}
